\section{Claw：Agent 的持久化层}

\subsection{Dobby the Elf Claw}

Karpathy 分享了他的 claw（持久化 agent）实验："Dobby the Elf Claw"是他为家庭自动化搭建的系统：

\begin{itemize}
    \item Agent 自动扫描局域网，发现 Sonos 音响系统并逆向工程其 API
    \item 同样的方式接管了灯光、HVAC、窗帘、泳池/SPA 控制
    \item 通过 Quinn 模型分析安全摄像头视频，检测变化后发送 WhatsApp 通知（如"FedEx 卡车刚停在门前"）
    \item 所有交互通过 WhatsApp 自然语言完成——"Dobby, sleepy time" 就关掉所有灯
\end{itemize}

\begin{knowledgebox}{从六个 App 到一个 WhatsApp}
Karpathy 原来使用六个不同的 app 管理智能家居，现在全部被 Dobby 替代。这指向一个更深层的问题：\textbf{App Store 中的大量 app 可能不应该存在}。它们应该只是 API endpoint，由 agent 作为"智能胶水"来调用和编排。
\end{knowledgebox}

\begin{figure}[H]
\centering
\includegraphics[width=0.82\textwidth]{frame-dobby-claw.jpg}
\caption{视频约 00:09:39：Dobby the Elf Claw 作为持久化 agent 的家庭自动化入口，体现了 memory、tool use 与 message channel 的结合}
\end{figure}

\subsection{从 Claw 到 Agent-First 世界}

\begin{importantbox}{软件行业的重构}
Agent 的客户不再是人类，而是其他 agent。行业需要从"human-first UI"重构为"agent-first API"。Karpathy 认为这种重构将是实质性的——treadmill、智能家居、各种 SaaS 工具都应该暴露 API 而不是构建 UI。
\end{importantbox}

\begin{knowledgebox}{Open Claw 的五个创新}
Karpathy 对 Peter 的 Open Claw 给予高度评价，认为其同时在五个方向创新：
\begin{enumerate}
    \item \textbf{Soul Document}——精心设计的 agent 人格（compelling personality）
    \item \textbf{Memory System}——比默认的 context compaction 更复杂的记忆系统
    \item \textbf{Persistence}——不依赖用户交互、持续运行的 loop
    \item \textbf{WhatsApp 入口}——统一的自然语言交互门户
    \item \textbf{人格调校}——如 Claude 的"praise calibration"让用户觉得表扬是真诚的
\end{enumerate}
\end{knowledgebox}


